\documentclass[11pt,a4paper]{article}
\usepackage[utf8]{inputenc}
\usepackage[spanish]{babel}
\usepackage{amsmath}
\usepackage{amsfonts}
\usepackage{amssymb}
\usepackage{geometry}
\usepackage{booktabs}
\usepackage{hyperref}
\usepackage{microtype}
\usepackage{listings}
\usepackage{xcolor}

\geometry{top=2.5cm, bottom=2.5cm, left=2.5cm, right=2.5cm}

\definecolor{codegray}{rgb}{0.5,0.5,0.5}
\definecolor{codegreen}{rgb}{0,0.6,0}
\definecolor{codeblue}{rgb}{0,0,0.8}
\definecolor{backcolour}{rgb}{0.95,0.95,0.92}

\listingsdefinelanguage{Python}{
  keywords={for, in, range, if, elif, else, break, continue, return, import, as, from},
  keywordstyle=\color{codeblue}\bfseries,
  ndkeywords={np, plt, tqdm},
  ndkeywordstyle=\color{purple}\bfseries,
  identifierstyle=\color{black},
  sensitive=false,
  comment=[l]{\#},
  commentstyle=\color{codegreen}\itshape,
  stringstyle=\color{red},
  morestring=[b]',
  morestring=[b]"
}

\lstset{
  language=Python,
  backgroundcolor=\color{backcolour},
  commentstyle=\color{codegreen},
  keywordstyle=\color{codeblue},
  numberstyle=\tiny\color{codegray},
  stringstyle=\color{red},
  basicstyle=\ttfamily\small,
  breakatwhitespace=false,
  breaklines=true,
  captionpos=b,
  keepspaces=true,
  numbers=left,
  numbersep=5pt,
  showspaces=false,
  showstringspaces=false,
  showtabs=false,
  tabsize=4
}

\title{\textbf{Definición de Variables, Parámetros y Dinámica Temporal en el Modelo de Dunas}}
\author{Documentación de Repositorio}
\date{\today}

\begin{document}

\maketitle

\section{Introducción}
Este documento detalla la interpretación física, matemática y el rol en la estabilidad numérica de los principales parámetros geométricos y dinámicos configurados en los scripts del directorio \texttt{03\_vortices}, específicamente enfocados en la migración de dunas, la segregación de sedimentos bidispersos y la resolución de la ecuación de Exner. Adicionalmente, se incluye una disección analítica del bucle principal de evolución temporal.

\section{Variables Geométricas de Frontera: \texttt{L\_flat\_left} y \texttt{L\_flat\_right}}
Estas variables definen la extensión de las zonas planas horizontales (lecho plano) antes y después del cuerpo de la duna principal dentro del canal de simulación.

\begin{itemize}
    \item \textbf{\texttt{L\_flat\_left}} ($0.05$ m / $5$ cm): Representa la longitud del tramo plano horizontal ubicado \textit{aguas arriba} (\textit{upstream}) de la duna, previo al inicio de la pendiente de barlovento (\textit{stoss}) en $x = 0$.
    \item \textbf{\texttt{L\_flat\_right}} ($0.50$ m / $50$ cm): Representa la longitud del tramo plano horizontal ubicado \textit{aguas abajo} (\textit{downstream}) de la duna, posterior al fin de la pendiente de sotavento (\textit{lee}) en $x = L_{\text{dune}}$.
\end{itemize}

\subsection{Efecto y Utilidad en la Simulación}
\begin{enumerate}
    \item \textbf{Definición del Dominio Espacial:} Establecen las coordenadas extremas de la grilla de cálculo mediante:
    \begin{align}
        x_{\text{min}} &= -L_{\text{flat\_left}} \\
        x_{\text{max}} &= L_{\text{dune}} + L_{\text{flat\_right}}
    \end{align}
    \item \textbf{Prevención de Colisión de Borde:} La velocidad de migración de la duna ($c$) depende de factores del transporte y no de estas variables. Sin embargo, al incrementarse el tiempo físico de simulación ($t_{\text{max}}$), un mayor valor de \texttt{L\_flat\_right} evita que la duna alcance el contorno derecho $x_{\text{max}}$, donde la condición de borde impone que la altura sea cero ($y_{\text{bed}} = 0$), evitando la deformación o destrucción artificial de la duna al chocar contra la frontera.
\end{enumerate}

\section{Parámetros de Depo-Distribución: $\lambda_s$ y $\lambda_l$}
Representan los coeficientes de decaimiento exponencial de deposición para las fracciones fina (\textit{small}) y gruesa (\textit{large}) de la mezcla de arena.

\begin{itemize}
    \item $\lambda_s = 60.0 \text{ m}^{-1}$: Tasa de decaimiento para granos finos ($d_s = 0.3 \text{ mm}$).
    \item $\lambda_l = 180.0 \text{ m}^{-1}$: Tasa de decaimiento para granos gruesos ($d_l = 1.0 \text{ mm}$).
\end{itemize}

\subsection{Relación con la Longitud de Viaje}
El inverso matemático de cada tasa de decaimiento, $\Lambda = 1/\lambda$, representa la \textbf{longitud característica de viaje} de las partículas suspendidas o en tránsito tras superar la cresta de la duna:
\begin{align}
    \Lambda_s &= \frac{1}{\lambda_s} = \frac{1}{60} \approx 0.0167 \text{ m} \approx 1.67 \text{ cm} \\
    \Lambda_l &= \frac{1}{\lambda_l} = \frac{1}{180} \approx 0.0056 \text{ m} \approx 0.56 \text{ cm}
\end{align}

\subsection{Mecanismo de Segregación Sedimentaria}
Dado que $\lambda_l > \lambda_s$, la longitud de asentamiento de los granos gruesos es significativamente menor que la de los finos. Al cruzar la cresta:
\begin{itemize}
    \item El sedimento grueso se deposita rápidamente en la zona superior de sotavento, cerca de la cresta.
    \item El sedimento fino viaja más lejos en suspensión, depositándose preferentemente en la parte inferior o pie de la pendiente.
\end{itemize}
Este comportamiento se modela a través del decaimiento del transporte en sotavento:
\begin{equation}
    q_s(x) = q_{s,\text{crest}} \left( \phi_{s,\text{bulk}} e^{-\lambda_s(x - x_{\text{crest}})} + \phi_{l,\text{bulk}} e^{-\lambda_l(x - x_{\text{crest}})} \right)
\end{equation}

\section{Coeficiente de Difusión por Pendiente: $D_{\text{grav}}$}
La variable \textbf{\texttt{D\_grav}} ($D_{\text{grav}} = 1.5 \times 10^{-7} \text{ m}^2/\text{s}$) es el coeficiente de difusividad morfológica incorporado en la ecuación de Exner:
\begin{equation}
    \frac{\partial y}{\partial t} = -\frac{1}{1 - p} \frac{\partial q_s}{\partial x} + D_{\text{grav}} \frac{\partial^2 y}{\partial x^2}
\end{equation}
donde $p$ es la porosidad del lecho y $q_s$ es la tasa de transporte longitudinal.

\subsection{Propósitos Clave}
\begin{enumerate}
    \item \textbf{Efecto Gravitacional Físico:} El término de curvatura $D_{\text{grav}} \frac{\partial^2 y}{\partial x^2}$ simula el transporte de sedimentos ladera abajo debido a la gravedad. Tiende a desgastar las cúspides (curvatura negativa) y rellenar las depresiones (curvatura positiva), suavizando el perfil.
    \item \textbf{Estabilidad Numérica (Condición de Von Neumann):} Para un esquema de diferencias finitas explícito en tiempo y centrado en espacio, la estabilidad del término difusivo exige que el número de Fourier ($Fo$) sea menor o igual a $0.5$:
    \begin{equation}
        Fo = \frac{D_{\text{grav}} \Delta t}{\Delta x^2} \le 0.5
    \end{equation}
    Bajo los parámetros de la malla ($\Delta x \approx 0.93\text{ mm}$ y $\Delta t = 2.0\text{ s}$), el valor calibrado arroja:
    \begin{equation}
        Fo = \frac{1.5 \times 10^{-7} \cdot 2.0}{(9.3 \times 10^{-4})^2} \approx 0.34
    \end{equation}
    lo que garantiza una simulación libre de oscilaciones espurias a escala de grilla.
\end{enumerate}

\section{Bucle Temporal de Simulación (\texttt{for step in range(Nt)})}
El bucle principal estructurado en la línea 91 corre paso a paso la evolución de la duna. Se divide en cinco etapas fundamentales en cada paso de tiempo $\Delta t$:

\subsection{Fase 1: Seguimiento Geométrico de la Cresta}
En primer lugar, se identifica la cresta como el máximo absoluto del perfil actual del lecho $y(x)$:
\begin{align}
    i_{\text{crest}} &= \operatorname{argmax}(y_{\text{bed}}) \\
    x_{\text{crest}} &= x_{\text{grid}}[i_{\text{crest}}] \\
    x_{\text{start}} &= x_{\text{crest}} - L_{\text{stoss}}
\end{align}
Esto permite definir dinámicamente el inicio de la duna a barlovento.

\subsection{Fase 2: Distribución del Flujo de Sedimentos ($q_s$)}
Se calcula vectorialmente el flujo local de transporte de sedimentos $q_s(x)$ a lo largo de la cuadrícula:
\begin{equation}
q_s(x) = \begin{cases} 
0, & \text{si } x \le x_{\text{start}} \quad \text{(Zona plana aguas arriba)} \\
q_{s,\text{crest}} \frac{1 - e^{-(x - x_{\text{start}})/L_{\text{sat}}}}{1 - e^{-L_{\text{stoss}}/L_{\text{sat}}}}, & \text{si } x_{\text{start}} < x \le x_{\text{crest}} \quad \text{(Barlovento / Crecimiento)} \\
q_{s,\text{crest}} \left( \phi_s e^{-\lambda_s(x - x_{\text{crest}})} + \phi_l e^{-\lambda_l(x - x_{\text{crest}})} \right), & \text{si } x > x_{\text{crest}} \quad \text{(Sotavento / Decaimiento)}
\end{cases}
\end{equation}

\subsection{Fase 3: Evolución del Lecho mediante Diferencias Finitas}
Se resuelve la ecuación de Exner utilizando diferencias finitas espaciales explícitas en el tiempo para actualizar la altura interna del lecho $y_i^{n+1}$:
\begin{equation}
    y_i^{n+1} = y_i^n - \left( \frac{\Delta t}{1-p} \right) \left( \frac{q_s[i] - q_s[i-1]}{\Delta x} \right) + D_{\text{grav}} \Delta t \left( \frac{y_{i+1}^n - 2y_i^n + y_{i-1}^n}{\Delta x^2} \right)
\end{equation}
con condiciones de frontera de lecho fijo en los extremos del dominio ($y_0 = 0.0$, $y_{Nx-1} = 0.0$).

\subsection{Fase 4: Modelo de Avalanchas Gravitatorias}
Si el lecho excede el ángulo de reposo estático crítico $\theta_{\text{static}}$ debido a la deposición en la cara de sotavento:
\begin{equation}
    \frac{y_i - y_{i+1}}{\Delta x} > \tan(\theta_{\text{static}} - \theta_{\text{stoss}})
\end{equation}
se transfiere una fracción del volumen de sedimento ladera abajo de forma iterativa hasta alcanzar el ángulo dinámico $\theta_{\text{dynamic}}$:
\begin{align}
    \Delta h &= \frac{1}{2} \left( \frac{y_i - y_{i+1}}{\Delta x} - \tan(\theta_{\text{dynamic}} - \theta_{\text{stoss}}) \right) \Delta x \\
    y_i &\leftarrow y_i - \Delta h \\
    y_{i+1} &\leftarrow y_{i+1} + \Delta h
\end{align}

\subsection{Fase 5: Registro de Estratigrafía Física y Vórtices}
Para cada celda de la malla, si hay deposición ($y_i^{n+1} > y_i^n$), se calcula la composición de finas $\phi_{\text{dep}}$ depositada en los índices verticales correspondientes de la cuadrícula de almacenamiento:
\begin{equation}
    \phi_{\text{dep}} = \phi_{\text{dep,base}} + \Delta\phi_{\text{vortex}}
\end{equation}
donde el efecto del vórtice (contracorriente en sotavento) modela el arrastre y convergencia de finas hacia la cresta dentro de la burbuja de separación de longitud $L_{\text{sep}}$:
\begin{equation}
\Delta\phi_{\text{vortex}} = \begin{cases}
0.15 \cos\left(\frac{\pi (x - x_{\text{crest}})}{L_{\text{sep}}}\right), & \text{si } x - x_{\text{crest}} \le L_{\text{sep}} \\
-0.15, & \text{si } x - x_{\text{crest}} > L_{\text{sep}}
\end{cases}
\end{equation}
Esta deposición también se modula periódicamente para reflejar avalanchas intermitentes, y el resultado es mapeado a la cuadrícula estratigráfica 2D.

\end{document}
